\documentclass[conference]{IEEEtran}
\IEEEoverridecommandlockouts
\usepackage{cite}
\usepackage{amsmath,amssymb,amsfonts}
\usepackage{graphicx}
\usepackage{textcomp}
\usepackage{xcolor}
\def\BibTeX{{\rm B\kern-.05em{\sc i\kern-.025em b}\kern-.08em
    T\kern-.1667em\lower.7ex\hbox{E}\kern-.125emX}}
\begin{document}

\title{Graph Neural Networks for Protein-Ligand Binding Affinity Prediction\\}

\author{
\IEEEauthorblockN{Kai Zuang}
\IEEEauthorblockA{\textit{Cornell University}\\
New York, USA \\
kz396@cornell.edu}
\and
\IEEEauthorblockN{Jiwon Jeong}
\IEEEauthorblockA{\textit{Cornell University}\\
New York, USA \\
jj835@cornell.edu}
}

\maketitle

\begin{abstract}
Predicting protein-ligand binding affinity is a central challenge in computational drug discovery, where experimental screening methods such as isothermal titration calorimetry and surface plasmon resonance are accurate but costly and low-throughput. In this work, we study how the choice of graph construction for the protein affects binding-affinity prediction when the neural-network architecture is held fixed. Our fixed backbone employs two parallel Graph Attention Network (GAT) encoders that independently embed the protein and ligand graphs, which are fused via cross-attention and passed to a regression head to produce a scalar affinity estimate. We vary the protein graph along two orthogonal axes---the residue-to-residue distance ($C_\alpha$--$C_\alpha$ distance vs.\ minimum heavy-atom distance) and the adjacency rule (direct Euclidean proximity vs.\ Isomap-style geodesic on a spectrally-embedded manifold)---at cutoffs of $3$ and $5$ \AA. We train and evaluate on LP-PDBBind, a leak-proof reorganization of PDBBind v2020 that controls for protein sequence and ligand structural similarity across splits, enabling rigorous generalization assessment to novel molecular structures. We report mean squared error and Pearson correlation on the held-out test set for each graph construction, isolating the contribution of representation choice from architectural choice and assessing which geometric inductive bias best supports generalization to novel protein-ligand complexes.
\end{abstract}

\section{Introduction}
Proteins are the dominant targets in modern drug discovery: 1,578 FDA-approved drugs act through 667 human-genome-derived proteins and 189 pathogen-derived proteins \cite{santos2017comprehensive}. In the drug discovery process, binding affinity is a core attribute, as it quantifies how strongly a protein interacts with small molecules, or ligands. Drug-target binding affinity serves as a vital indicator of how well a drug interacts with its target, essential for achieving intended therapeutic outcomes \cite{shaker2021drug}. Binding affinity screening is therefore a standard and critical step in early-stage drug discovery pipelines. To illustrate its role as a filter: the screening attrition rate in current drug discovery protocols suggests that only one marketable drug emerges from approximately one million screened compounds \cite{berdasco2020hts}, and binding affinity is among the primary criteria used to eliminate candidates at this stage. Experimentally, binding affinities are most commonly measured using isothermal titration calorimetry (ITC) and surface plasmon resonance (SPR). ITC follows the heat change when a test compound binds to a target protein, enabling precise measurement of affinity without requiring competing molecules or molecular labels \cite{freyer2008itc}. However, ITC and SPR are lower-throughput techniques \cite{freyer2008itc}, and both methods are costly and difficult to scale for \textit{in silico} and mass drug discovery pipelines. In this project, we aim to use graph data science to predict binding affinities by leveraging graph neural networks, graphical models, and different graph representations for three-dimensional molecular structures. We aim to generalize to unseen structures so that the model can be applied in drug discovery pipelines, where many candidates are novel.

To perform this prediction task, we use the LP-PDBBind dataset, a reorganized and leak-proof version of PDBBind v2020. LP-PDBBind carefully prepares a cleaned set of non-covalent protein-ligand binders split into training, validation, and test sets to control for data leakage, defined as proteins and ligands with high sequence and structural similarity \cite{ha2023lp}, drawing from the 14,108 protein-ligand complexes in the PDBBind v2020 general set \cite{liu2015pdbbind}. To robustly train models without data leakage, we adopt LP-PDBBind's similarity-based train-test split, controlling for both protein sequence similarity and ligand structural similarity, such that the validation and test sets contain molecularly distinct structures \cite{ha2023lp}.

\section{Problem Setup}
The central question of this work is how the choice of graph construction for the protein affects binding-affinity prediction, with the neural-network architecture, training procedure, and ligand representation all held fixed. The PDBBind dataset provides experimentally determined protein-ligand complex structures in PDB format, giving direct access to 3D atomic coordinates from which the protein graph can be built in several principled ways.

We vary the protein graph along two orthogonal axes. The first axis is the residue-to-residue distance $d(i,j)$: either the $C_\alpha$--$C_\alpha$ distance $d_{\mathrm{CA}}(i,j) = \|\mathbf{x}_{C_\alpha^i} - \mathbf{x}_{C_\alpha^j}\|_2$, which treats each residue as a single point and inherits the Euclidean metric of $\mathbb{R}^3$, or the minimum heavy-atom distance $d_{\mathrm{MHAD}}(i,j) = \min_{a \in \mathcal{A}_i, b \in \mathcal{A}_j} \|\mathbf{x}_a - \mathbf{x}_b\|_2$, which treats each residue as a set of atoms and picks the closest pair. The second axis is how $d$ is converted into an unweighted adjacency: a direct Euclidean proximity rule $A_{ij} = \mathbf{1}[d(i,j) \leq \tau]$, or an Isomap-style geodesic rule where edges link residues whose shortest path on the neighborhood graph of radius $r = \tau$ has length at most $\tau$. In parallel we record a normalized-Laplacian spectral embedding of the same neighborhood graph to characterize the manifold on which the geodesic is computed. Combined with two cutoffs $\tau \in \{3, 5\}$ \AA, these axes define the grid of protein-graph constructions we compare empirically in Sec.~\ref{sec:residue_graphs}.

The ligand graph is held fixed across all experiments: heavy atoms as nodes and covalent bonds as edges, with node features $\mathbf{h}_j$ encoding atomic number, hybridization, and formal charge, and edge features $\mathbf{e}_{jk}$ encoding bond type. Protein node features $\mathbf{h}_i$ encode amino-acid type, secondary-structure label, and solvent-accessible surface area, and protein edge features $\mathbf{e}_{ij}$ encode the underlying residue distance $d(i,j)$. Once the graphs are constructed, they are passed into separate GAT encoders, pooled to graph-level representations $\mathbf{z}_p$ and $\mathbf{z}_\ell$, and fused via cross-attention before a regression head predicts the binding affinity $\Delta G$. The model architecture is described in Sec.~\ref{AA} and is held fixed across all protein-graph constructions so that any differences in test-set performance can be attributed to the graph-representation choice.

\subsection{Solution}\label{AA}
Our proposed solution is a Graph Neural Network (GNN) that operates over the full molecular graph structure of both the protein and ligand, learning rich structural representations by aggregating information across nodes and edges. Given a protein-ligand pair $(G_p, G_\ell)$, the model $f_\theta(G_p, G_\ell) \to \mathbb{R}$ predicts the binding affinity $\Delta G$ end-to-end, leveraging the complete 3D topology of each molecule rather than relying on hand-crafted descriptors or partial structure summaries.

The architecture consists of two parallel Graph Attention Network (GAT) encoders, one for each modality. The protein encoder takes $G_p = (\mathcal{V}_p, \mathcal{E}_p)$ where each node $v_i \in \mathcal{V}_p$ represents an amino acid residue with feature vector $\mathbf{h}_i^{(0)} \in \mathbb{R}^{d_p}$ encoding residue type, secondary structure, and solvent accessibility, and each edge $e_{ij} \in \mathcal{E}_p$ encodes bond type or spatial distance $\|\mathbf{x}_i - \mathbf{x}_j\|_2$. The ligand encoder takes $G_\ell = (\mathcal{V}_\ell, \mathcal{E}_\ell)$ where nodes represent heavy atoms with features $\mathbf{h}_j^{(0)} \in \mathbb{R}^{d_\ell}$ encoding atomic number, hybridization, and formal charge, and edges encode covalent bond type. Each encoder applies $L$ layers of multi-head graph attention, updating node representations as
\begin{equation}
    \mathbf{h}_i^{(l+1)} = \sigma\!\left(\sum_{j \in \mathcal{N}(i)} \alpha_{ij}^{(l)}\, \mathbf{W}^{(l)} \mathbf{h}_j^{(l)}\right)
\end{equation}
where $\alpha_{ij}^{(l)}$ are the learned attention coefficients and $\mathbf{W}^{(l)}$ is a trainable weight matrix, with residual connections and layer normalization applied between each layer. Node embeddings are then aggregated via global mean pooling to obtain fixed-size graph-level representations $\mathbf{z}_p, \mathbf{z}_\ell \in \mathbb{R}^{d_h}$.

The two graph-level representations are fused via a cross-attention layer in which $\mathbf{z}_p$ and $\mathbf{z}_\ell$ serve as mutual queries and keys, producing an interaction-aware vector $\mathbf{z}_{\text{fused}} = \text{CrossAttn}(\mathbf{z}_p, \mathbf{z}_\ell)$. This fused representation is passed to an MLP regression head
\begin{equation}
    \hat{y} = \mathbf{W}_2\, \sigma(\mathbf{W}_1\, \mathbf{z}_{\text{fused}} + \mathbf{b}_1) + \mathbf{b}_2
\end{equation}
which outputs the predicted binding affinity $\hat{y} \in \mathbb{R}$ in units of $\text{kcal/mol}$. The model is trained end-to-end by minimizing the mean squared error loss
\begin{equation}
    \mathcal{L} = \frac{1}{N}\sum_{k=1}^{N}\left(\hat{y}_k - y_k\right)^2
\end{equation}
over $N$ training complexes, using the AdamW optimizer with weight decay $\lambda$ and a cosine annealing learning rate schedule.

\section{Analysis and Numerical Results}
As an initial feasibility check, we characterized the LP-PDBBind splits to confirm that the dataset supports our modeling setup. The cleaned dataset contains 19{,}443 protein-ligand complexes partitioned into train (11{,}513), validation (2{,}422), and test (4{,}860) sets, with the remainder unassigned. Complexes span 13 protein classes dominated by hydrolases (5{,}733) and transferases (5{,}291), confirming broad enzyme-class coverage for the GAT encoders to learn over. Table~\ref{tab:split_stats} summarizes average protein length, ligand heavy-atom count, and binding affinity (reported as $-\log_{10}\mathrm{K_d/K_i}$) per split. Train complexes feature longer proteins (402 residues) than validation/test (341--342), while ligand size is comparable in train and validation ($\sim$34 heavy atoms) but smaller in test (28). Mean binding affinity shifts from 6.21 (train) to 6.58 (test) to 6.88 (validation), indicating the held-out splits are enriched for tighter binders --- a useful stress test for generalization to novel structures.

\begin{table}[htbp]
\caption{LP-PDBBind dataset statistics by split.}
\label{tab:split_stats}
\begin{center}
\begin{tabular}{|l|c|c|c|c|}
\hline
\textbf{Split} & \textbf{Count} & \textbf{Avg. prot. len.} & \textbf{Avg. lig. atoms} & \textbf{Avg. affinity} \\
\hline
train & 11{,}513 & 402.1 & 35.1 & 6.21 \\
val   & 2{,}422  & 341.6 & 34.1 & 6.88 \\
test  & 4{,}860  & 339.8 & 28.4 & 6.58 \\
\hline
\end{tabular}
\end{center}
\end{table}

\subsection{Residue-Graph Constructions}\label{sec:residue_graphs}
Before training, we compare two principled axes of choice when converting a PDB structure into a residue-level graph, using PDB entry 7EJK (257 standard residues) as a case study. Nodes are residues; we consider the $0/1$ unweighted adjacency $A_{ij} = \mathbf{1}[\,d(i,j) \leq \tau\,]$ under a $2\times 2\times 2$ grid of choices: (i) the residue-to-residue distance $d$, (ii) whether that distance is used directly (Euclidean proximity) or lifted through a manifold (geodesic), and (iii) the cutoff $\tau \in \{3, 5\}$ \AA.

\paragraph{Two residue distances} We consider two definitions of $d(i,j)$. The first is the $C_\alpha$--$C_\alpha$ distance, $d_{\mathrm{CA}}(i,j) = \|\mathbf{x}_{C_\alpha^i} - \mathbf{x}_{C_\alpha^j}\|_2$, which treats each residue as a single point in $\mathbb{R}^3$ and inherits the usual Euclidean metric. The second is the minimum heavy-atom distance $d_{\mathrm{MHAD}}(i,j) = \min_{a \in \mathcal{A}_i,\, b \in \mathcal{A}_j} \|\mathbf{x}_a - \mathbf{x}_b\|_2$, where $\mathcal{A}_i$ is the set of heavy (non-hydrogen) atoms of residue $i$. MHAD treats each residue as a set of atoms and picks the closest pair, which makes the small $3$--$5$ \AA\ cutoff regime meaningful: consecutive $C_\alpha$ atoms are roughly $3.8$ \AA\ apart, so $\tau = 3$ \AA\ on $C_\alpha$ distance yields an essentially empty graph. Note also that $d_{\mathrm{MHAD}}$ is a minimum over pairs and \emph{does not} satisfy the triangle inequality, so it is not a metric; $d_{\mathrm{CA}}$ is.

\paragraph{Euclidean vs.\ geodesic adjacency} For either distance we build both a direct adjacency $A_{ij} = \mathbf{1}[d(i,j) \leq \tau]$ and a geodesic adjacency constructed in the Isomap style. The geodesic construction first forms a neighborhood graph $G_r$ whose edges link residues with $d(i,j) \leq r$ and carry weight $d(i,j)$. The geodesic distance $d_{\mathrm{geo}}^{(r)}(i,j)$ is the shortest path in $G_r$, which approximates the intrinsic metric on the manifold supported by the residue coordinates. In parallel, a normalized-Laplacian spectral embedding of $G_r$ with a Gaussian affinity $W_{ij} = \exp(-d(i,j)^2 / 2\sigma^2)$, $\sigma = r/2$, gives a low-dimensional coordinate system on the same manifold. The adjacency uses $\tau = r$ in each case: $(r,\tau) \in \{(3,3), (5,5)\}$ \AA.

\begin{table}[htbp]
\caption{Residue-graph adjacency statistics for 7EJK ($n = 257$, $\binom{n}{2} = 32{,}896$ possible edges). All eight configurations use a cutoff $\tau$ equal to the neighborhood radius $r$.}
\label{tab:graph_constructions}
\begin{center}
\begin{tabular}{|l|l|c|c|c|c|}
\hline
\textbf{Dist.} & \textbf{Adj.} & $\boldsymbol{\tau}$ (\AA) & \textbf{Edges} & \textbf{Density} & \textbf{Seq.\,/\,Non-seq.} \\
\hline
$C_\alpha$ & Euclid.  & 3 & 2     & 0.0001 & 2\,/\,0 \\
$C_\alpha$ & Geodesic & 3 & 2     & 0.0001 & 2\,/\,0 \\
$C_\alpha$ & Euclid.  & 5 & 293   & 0.0089 & 254\,/\,39 \\
$C_\alpha$ & Geodesic & 5 & 293   & 0.0089 & 254\,/\,39 \\
MHAD       & Euclid.  & 3 & 457   & 0.0139 & 254\,/\,203 \\
MHAD       & Geodesic & 3 & 691   & 0.0210 & 254\,/\,437 \\
MHAD       & Euclid.  & 5 & 1{,}299 & 0.0395 & 254\,/\,1{,}045 \\
MHAD       & Geodesic & 5 & 1{,}749 & 0.0532 & 254\,/\,1{,}495 \\
\hline
\end{tabular}
\end{center}
\end{table}

Table~\ref{tab:graph_constructions} reports all eight configurations. Three observations stand out.

First, the $C_\alpha$ graph at $\tau = 3$ \AA\ contains only $2$ edges, both backbone-consecutive pairs with a shortened $C_\alpha$--$C_\alpha$ separation (min $2.93$ \AA, likely $cis$-peptide geometry). This confirms that a naive point-per-residue representation cannot support the tight cutoffs used for contact maps and motivates MHAD.

Second, every configuration recovers $254$ of the $256$ backbone-consecutive pairs (the two missing ones correspond to chain breaks in the deposited structure), except the $C_\alpha$-$3$ \AA\ settings where only $2$ of these sequential edges fit under the cutoff. The $C_\alpha$-$5$ \AA\ graph then adds only $39$ non-sequential contacts, whereas MHAD at the same cutoff adds $1{,}045$, illustrating how much richer the contact set becomes once the full heavy-atom sidechain geometry is considered.

Third, and most informatively, the Euclidean and geodesic adjacencies coincide exactly for the $C_\alpha$ distance but diverge for MHAD. Under $C_\alpha$, $d_{\mathrm{CA}}$ is a proper metric in $\mathbb{R}^3$, so shortest-path distance on $G_r$ never improves on the direct edge and $A^{\mathrm{geo}} = A^{\mathrm{euc}}$. Under MHAD, $d_{\mathrm{MHAD}}$ violates the triangle inequality because it is a minimum over atom pairs, so a two-hop path $i \to k \to j$ through an intermediate residue can have total weight smaller than the direct $d_{\mathrm{MHAD}}(i,j)$. The geodesic construction repairs the metric property and in the process adds $234$ extra edges at $\tau = 3$ \AA\ and $450$ at $\tau = 5$ \AA. We interpret these bonus edges as second-shell contacts mediated through a shared neighbor residue, which are physically relevant for cooperative packing and allosteric signalling. The four nontrivial constructions (MHAD-Euclid and MHAD-geodesic at $\tau \in \{3, 5\}$ \AA) therefore define genuinely different inductive biases for the GAT encoder; we treat the MHAD-geodesic adjacency at $\tau = 5$ \AA\ as our default and will ablate the choice in subsequent experiments.

\begin{thebibliography}{00}

\bibitem{santos2017comprehensive}
R. Santos, O. Ursu, A. Gaulton, A. P. Bento, R. S. Donadi, C. G. Bologa, A. Karlsson, B. Al-Lazikani, A. Hersey, T. I. Oprea, and J. P. Overington,
``A comprehensive map of molecular drug targets,''
\textit{Nature Reviews Drug Discovery}, vol. 16, no. 1, pp. 19--34, 2017.

\bibitem{shaker2021drug}
B. Shaker, M.-S. Yu, J. S. Song, S. Ahn, J.-Y. Ryu, K.-S. Oh, and D. Na,
``Advancing drug safety in drug development: bridging computational predictions for enhanced toxicity prediction,''
\textit{Chemical Research in Toxicology}, vol. 34, no. 8, pp. 1914--1932, 2021.

\bibitem{berdasco2020hts}
E. P{\'e}rez-Herrero and A. Fern{\'a}ndez-Medarde,
``High throughput screening in drug discovery,''
\textit{European Journal of Pharmaceutics and Biopharmaceutics}, vol. 137, pp. 1--12, 2019.

\bibitem{freyer2008itc}
M. W. Freyer and E. A. Lewis,
``Isothermal titration calorimetry: experimental design, data analysis, and probing macromolecule/ligand binding and kinetic interactions,''
\textit{Methods in Cell Biology}, vol. 84, pp. 79--113, 2008.

\bibitem{ha2023lp}
E. J. Ha, C. T. Lwin, and J. D. Durrant,
``Leak proof {PDBBind}: a reorganized dataset of protein-ligand complexes for more generalizable binding affinity prediction,''
\textit{arXiv preprint arXiv:2308.09639}, 2023.

\bibitem{liu2015pdbbind}
Z. Liu, Y. Li, L. Han, J. Li, J. Liu, Z. Zhao, W. Nie, Y. Liu, and R. Wang,
``{PDB}-wide collection of binding data: current status of the {PDBbind} database,''
\textit{Bioinformatics}, vol. 31, no. 3, pp. 405--412, 2015.

\end{thebibliography}

\end{document}